\documentclass[11pt]{article}
\usepackage[margin=1in]{geometry}
\usepackage{amsmath,amssymb,amsthm}
\usepackage{algorithm}
\usepackage{algpseudocode}
\usepackage{booktabs}
\usepackage{hyperref}
\usepackage{tikz}
\usepackage{xcolor}
\usetikzlibrary{arrows.meta,positioning,fit,backgrounds,shapes.geometric,calc}

\newtheorem{theorem}{Theorem}
\newtheorem{proposition}{Proposition}
\newtheorem{lemma}{Lemma}
\newtheorem{definition}{Definition}

\newcommand{\lo}{\underline}
\newcommand{\hi}{\overline}
\newcommand{\indep}{\perp\!\!\!\perp}
\newcommand{\DSE}{\mathcal{D}_{\mathrm{SE}}}
\newcommand{\DLCN}{\mathcal{D}_{\mathrm{LCN}}}
\newcommand{\ALP}{\texttt{ApproxLP}}

\title{Exactness of ApproxLP on Chain-Graph LCNs\\
\large Why ApproxLP is an \emph{inner} bound of the strong extension, when the
inner bound is nonetheless tight, and how it pairs with Credal~VE to bracket the
LCN truth}
\author{IBM Research}
\date{}

\begin{document}
\maketitle

\begin{abstract}
ApproxLP (\texttt{lcn/inference/marginal/cn/approxlp.py}, class \texttt{ApproxLP}) is the
coordinate-descent member of the credal-network family: it reformulates credal marginal inference as
a multilinear program and solves it by iterated linearization --- fixing all local distributions but
one and picking that family's best extreme point, sweeping to a local optimum (Antonucci, de~Campos,
Huber \& Zaffalon, 2015). This note pins down three things about ApproxLP against the implementation.
\textbf{(i)~What it bounds.} Like every credal-network engine, ApproxLP works over the \emph{strong
extension} $\DSE$ of the compiled network, not the LCN's distribution set $\DLCN$. Uniquely in the
family, it is an \emph{inner} approximation: because coordinate descent visits only \emph{feasible}
product joints and converges to a \emph{local} optimum of a nonconvex multilinear program, the
returned interval sits \emph{inside} the strong-extension interval $[\lo q_{\mathrm{SE}},\hi
q_{\mathrm{SE}}]$. It is therefore \emph{not} a valid outer bound on $\DLCN$; it can be too tight.
\textbf{(ii)~When the inner bound is tight.} When every per-variable linearization has its global
optimum reachable by coordinate descent --- which holds on the singly-connected, sentence-free
instances where $\DSE=\DLCN$ and the multilinear objective is benign --- ApproxLP's inner interval
\emph{coincides} with the exact bound. We verify this on \texttt{d4\_biting}: ApproxLP returns
$P(A){=}[0.40,0.60]$, $P(B){=}[0.04,0.72]$, $P(C){=}[0.08,0.80]$, matching Credal~VE and the exact
oracle to four decimals --- a sound, tight inner bound (\texttt{IntervalBP}, after its two
message-passing bugs were fixed, now matches these too; \texttt{docs/ibp\_exactness.tex}). The inner
gap becomes strict only when coordinate descent stalls at a non-global local optimum.
\textbf{(iii)~Which
tightening schemes apply.} D1/D2/D3 are \emph{build-stage} knobs that ApproxLP inherits for
\emph{free} (tighter vertices $\Rightarrow$ tighter LPs). D4 (\texttt{coupling="cross-family"}) \emph{is}
implemented, as extra LP rows that keep every visited joint feasible --- but it is \emph{experimental
and not a certified} bound: when the center initialization already violates a constraint, descent can
stall at a degenerate point. D5 is the separate exact engine \texttt{CredalJT}. The natural use of
ApproxLP is as the \emph{inner} half of a two-sided bracket: ApproxLP (inner) with Credal~VE (outer)
on the \emph{same} vertices sandwiches the LCN truth.
\end{abstract}

\tableofcontents

\section{Background: the ApproxLP algorithm}\label{sec:bg}

\subsection{LCNs, the credal network, and the strong extension}
An LCN over binary atoms $\mathbf V=\{V_1,\dots,V_n\}$ is a set of probability \emph{sentences},
each Type-1 $\ell\le P(\varphi)\le u$ or Type-2 $\ell\le P(\varphi\mid\psi)\le u$, over
propositional formulas. With the conditional independencies of the \emph{Local Markov Condition}
(LMC) of the primal graph, they define a set $\DLCN\subseteq\Delta(2^{\mathbf V})$ of joints, and the
marginal task for a query atom $a$ is $\lo P(a)=\min_{p\in\DLCN}P(a)$, $\hi P(a)=\max_{p\in\DLCN}
P(a)$.

The \texttt{cn/} pipeline compiles the chain-graph LCN into a \emph{credal network}: a directed
graph over the families $\mathrm{ch}_v\mid\mathrm{pa}_v$ with interval-valued local credal sets
$K_v$. \texttt{CredalNetworkVertices.from\_lcn} enumerates the extreme points of every $K_v$
(\texttt{cnv.extreme\_points}); all engines --- Credal~VE, CCTE, IBP, ApproxLP, Credal~IJGP ---
consume that object. What they bracket is the \emph{strong extension}
\begin{equation}\label{eq:se}
  \DSE=\Bigl\{\textstyle\prod_v P_v \;:\; P_v\in K_v \text{ chosen independently per parent-config}\Bigr\},
\end{equation}
and, because the chain-graph factorization reproduces the LMC exactly,
\begin{equation}\label{eq:incl}
  \DLCN\;\subseteq\;\DSE ,
\end{equation}
with equality unless a \emph{cross-family sentence} couples choices the free product decouples.
Credal~VE returns \emph{exactly} $[\lo q_{\mathrm{SE}},\hi q_{\mathrm{SE}}]$
(\texttt{docs/cve\_exactness.tex}); ApproxLP approximates the \emph{same} target \emph{from the
inside}.

\subsection{Multilinear program and coordinate descent}\label{sec:cd}
Credal marginal inference over $\DSE$ is a \emph{multilinear program}: the objective
$P(\text{query}\mid\text{evidence})$ is a ratio that is linear in each family's distribution
$P_v$ when the others are fixed, and the feasible region is the product of the polytopes $K_v$
(\texttt{approxlp.py:41--49}). ApproxLP (\texttt{approxlp.py:306--380}, \texttt{\_coordinate\_descent})
solves it by iterated linearization:
\begin{enumerate}
  \item \textbf{Initialize} every local $P_v$ to the \emph{center} (mean of its extreme points)
    (\texttt{approxlp.py:112--131, 318}).
  \item \textbf{Sweep} (\texttt{approxlp.py:332--370}): for each family $v$ and each parent
    configuration, hold all other families fixed; the objective is then linear in $P_v$, so its
    optimum over the polytope $K_v$ is at a \emph{vertex}. Enumerate the extreme points, evaluate
    the objective for each by an exact variable-elimination pass on the resulting precise Bayes net
    (\texttt{approxlp.py:133--234}), and \emph{greedily commit} the best (\texttt{approxlp.py:359--370}).
  \item \textbf{Converge} (\texttt{approxlp.py:375--378}): stop when a full sweep improves nothing, or
    after \texttt{n\_iters} (default $50$).
\end{enumerate}
The run is repeated with \texttt{sense="min"} and \texttt{"max"} to get the lower and upper bound of
each query. The public API is
\texttt{run(evidence, n\_iters=50, coupling="off", verbosity=1)} (\texttt{approxlp.py:386--389}).

\section{What ApproxLP bounds: inner over the strong extension}\label{sec:bounds}

\begin{proposition}[ApproxLP is an inner bound]\label{prop:inner}
With \texttt{coupling="off"}, ApproxLP returns per-atom intervals \emph{contained in} the
strong-extension interval:
$[\lo q_{\mathrm{ALP}},\hi q_{\mathrm{ALP}}]\subseteq[\lo q_{\mathrm{SE}},\hi q_{\mathrm{SE}}]$.
In general it is therefore \emph{not} a valid outer bound on $\DLCN$.
\end{proposition}

\begin{proof}[Proof sketch]
Every state ApproxLP evaluates is a \emph{single feasible product joint}: at every point of the
sweep each family sits at one extreme point of its $K_v$ (or at the center, itself a convex
combination of them), so the assembled joint is a point of $\DSE$ (\eqref{eq:se}). The objective is
evaluated exactly on that precise Bayes net (\texttt{approxlp.py:133--234}). Hence every value
ApproxLP ever records is $P(a)$ for some $P\in\DSE$, and lies in $[\lo q_{\mathrm{SE}},\hi
q_{\mathrm{SE}}]$. The reported lower (resp.\ upper) bound is the smallest (largest) such value found
along the descent trajectory, which is a $\min$ (resp.\ $\max$) over a \emph{subset} of $\DSE$ and
therefore no smaller (no larger) than the true $\DSE$ extremum. This is the exact opposite of
Credal~VE, which enumerates \emph{all} vertex combinations and so attains the global $\DSE$
extremum: $\lo q_{\mathrm{SE}}\le\lo q_{\mathrm{ALP}}\le\hi q_{\mathrm{ALP}}\le\hi q_{\mathrm{SE}}$.
\end{proof}

\noindent The inner gap is strict exactly when coordinate descent \emph{stalls}: the multilinear
program is nonconvex, so a greedy sweep from the center can reach a local optimum that is not the
global $\DSE$ extremum (no single-family move improves, yet a coordinated multi-family jump would).
The inner gap is \emph{zero} when the reachable local optimum is global --- the tightness question of
Section~\ref{sec:tight}.

\section{When the inner bound is tight}\label{sec:tight}

By Proposition~\ref{prop:inner}, ApproxLP equals the exact LCN bound when \emph{both} (a) its inner
interval coincides with $[\lo q_{\mathrm{SE}},\hi q_{\mathrm{SE}}]$ (coordinate descent reaches the
global $\DSE$ extremum) \emph{and} (b) $\DSE=\DLCN$ (singly connected, no cross-family sentence, as
for Credal~VE). Condition (b) is the same geometric condition analyzed in
\texttt{docs/cve\_exactness.tex}; condition (a) is an optimization-landscape property of the
multilinear program.

\begin{proposition}[Tightness on benign instances]\label{prop:tight}
On a clean singly-connected, cross-family-sentence-free LCN whose per-family linearizations each
expose their global optimum to a greedy sweep from the center, ApproxLP's inner interval coincides
with the exact LCN marginal bounds.
\end{proposition}

\begin{proof}[Proof sketch]
(b) gives $\DSE=\DLCN$ exactly as in \texttt{docs/cve\_exactness.tex} (a chain-graph product satisfies
every LMC; with no cross-family sentence the free product cannot leave $\DLCN$). Under (a), the
descent reaches the global $\DSE$ extremum, so by Proposition~\ref{prop:inner}'s chain of
inequalities the inner interval equals $[\lo q_{\mathrm{SE}},\hi q_{\mathrm{SE}}]=[\lo
q_{\mathrm{LCN}},\hi q_{\mathrm{LCN}}]$.
\end{proof}

\subsection{Empirical confirmation on \texttt{d4\_biting}}\label{sec:empirical}
\texttt{examples/d4\_biting.lcn} (three atoms $A\to B$, $A\to C$; families $P(A),P(B\mid A),P(C\mid A)$;
one cross-family sentence $P(B\land C)\le0.05$) is the family's standard stress instance. The exact
unconditional bounds are $P(A)=[0.40,0.60]$, $P(B)=[0.04,0.72]$, $P(C)=[0.08,0.80]$ (four independent
oracles agree; \texttt{docs/ccte\_exactness.tex}). Measured, with \texttt{coupling="off"}:
\[
  \text{\ALP: }P(A){=}[0.40,0.60],\quad P(B){=}[0.04,0.72],\quad P(C){=}[0.08,0.80],
\]
matching the exact bounds to four decimals. The engines on the \emph{same} vertices:
\begin{center}\small
\begin{tabular}{lccc}
\toprule
Engine on \texttt{d4\_biting} & $P(A)$ & $P(B)$ & $P(C)$ \\
\midrule
Exact $=$ Credal~VE $=$ oracle          & $[0.40,0.60]$ & $\mathbf{[0.04,0.72]}$ & $[0.08,0.80]$ \\
\textbf{ApproxLP} (\texttt{coupling="off"})      & $[0.40,0.60]$ & $\mathbf{[0.04,0.72]}$ & $[0.08,0.80]$ \\
\texttt{IntervalBP} (\texttt{interval}, fixed) & $[0.40,0.60]$ & $\mathbf{[0.04,0.72]}$ & $[0.08,0.80]$ \\
CCTE (\texttt{coupling="off"})          & $[0.40,0.60]$ & $\mathbf{[0.05,0.65]}$ & $[0.08,0.80]$ \\
\bottomrule
\end{tabular}
\end{center}
On this instance ApproxLP's inner interval reaches the exact bound, and \texttt{IntervalBP} matches
it too now that its two message-passing bugs are fixed (before the fix it returned the inner
$[0.5,0.5]$ / $[0.05,0.65]$ / $[0.10,0.75]$; \texttt{docs/ibp\_exactness.tex}). CCTE remains inner on
$P(B)$ (\texttt{docs/ccte\_exactness.tex}). The reason ApproxLP succeeds here is that, for these
unconditional marginals, the objective's global $\DSE$ extremum \emph{is} a coordinate-descent fixed
point reachable from the center --- condition (a) holds. On harder (nonconvex, more coupled)
instances condition (a) can fail and the inner gap becomes strict; then ApproxLP under-reports the
width and must be read as an inner bound only.

\subsection{Topology by topology}
\paragraph{Markov chain, tree, polytree.} Singly connected, no cross-family sentence, so
$\DSE=\DLCN$; the per-family linearizations are simple and coordinate descent reaches the global
extremum, so ApproxLP is tight (equal to \texttt{ExactInference} / Credal~VE) on
\texttt{examples/chain.lcn}, \texttt{tree.lcn}, \texttt{polytree.lcn}. Verified on a two-atom chain
$x_0\!\to\!x_1$: ApproxLP returns $P(x_0){=}[0.300,0.500]$, $P(x_1){=}[0.400,0.680]$, matching
Credal~VE exactly (the fixed \texttt{IntervalBP} now agrees on this instance too;
\texttt{docs/ibp\_exactness.tex}).

\paragraph{Chain graph with compound nodes (\texttt{alarm}).} Still no cross-family sentence, so
$\DSE=\DLCN$; ApproxLP is tight provided descent does not stall on the larger compound-node
linearizations. This is the regime where pairing with Credal~VE (Section~\ref{sec:pairing}) is most
valuable as a stall detector.

\paragraph{Loopy DAG / cross-family sentence.} $\DSE\supsetneq\DLCN$ (or descent stalls). ApproxLP
remains an \emph{inner} bound of $\DSE$; it never over-covers, so it cannot certify an outer bound on
$\DLCN$. Use it as the inner half of a bracket, and use D5 (\texttt{CredalJT}) for a guaranteed exact
bound.

\section{Pairing with Credal~VE: a two-sided bracket}\label{sec:pairing}

Because ApproxLP is \emph{inner} over $\DSE$ and Credal~VE is \emph{exact} over $\DSE$, running both
on the \emph{same} enumerated vertices yields
\[
  \underbrace{[\lo q_{\mathrm{ALP}},\hi q_{\mathrm{ALP}}]}_{\text{inner}}
  \;\subseteq\;
  \underbrace{[\lo q_{\mathrm{SE}},\hi q_{\mathrm{SE}}]}_{\text{Credal~VE, exact over }\DSE}
  \;\supseteq\;
  \underbrace{[\lo q_{\mathrm{LCN}},\hi q_{\mathrm{LCN}}]}_{\text{truth, when }\DSE=\DLCN}.
\]
When the two intervals \emph{coincide}, coordinate descent reached the global $\DSE$ extremum and the
common value is certified as the exact $\DSE$ bound; when they \emph{differ}, the gap between them
is precisely the coordinate-descent stall, flagging that ApproxLP under-reported the width. This is
the recommended diagnostic use of ApproxLP (\texttt{docs/tighter\_approximation.tex}): the inner
ApproxLP and outer Credal~VE bracket the LCN truth from both sides on the tightened vertices.

\section{The tightening schemes (D1--D5) for ApproxLP}\label{sec:schemes}

The companion note \texttt{docs/tighter\_approximation.tex} identifies \textbf{Gap~A} (local-set
widening) and \textbf{Gap~B} (strong-extension widening) and five schemes.

\paragraph{Class 1 --- input-level (D1, D2, D3): free build-stage knobs.} These change only what
\texttt{CredalNetworkVertices.from\_lcn} produces; ApproxLP consumes tighter vertices without any code
change, and its per-variable LP already optimizes over the resulting (smaller) vertex polytope, so
its inner interval tightens toward $\DLCN$ from the inside. \textbf{D1} (\texttt{method="linear-tight"})
enriches each family LP with in-scope sentences and scope-restricted LMC equalities; \textbf{D2}
(\texttt{merge\_budget=$k$}) merges adjacent families; \textbf{D3} (\texttt{solver="scip"}) certifies
each family's local set globally.

\paragraph{Class 2 --- propagation-level coupling (D4, D5).}
\begin{itemize}
  \item \textbf{D4} (\texttt{coupling="cross-family"} on \texttt{ApproxLP.run}): implemented
    (\texttt{approxlp.py:265--304, 346--350}) by rejecting any vertex pick whose fully reconstructed
    joint violates a cross-family constraint --- i.e.\ adding D4's coupling rows to the per-variable LP,
    which keeps the inner-bound guarantee while shrinking the polytope
    (\texttt{docs/tighter\_approximation.tex}). \textbf{But it is not a certified bound}: the
    \texttt{run} docstring warns that ApproxLP ``is an inner, greedy coordinate-descent method and is
    the loosest fit for D4''; when the center initialization already violates a constraint, descent
    can find no feasible improving move and \emph{stall at a degenerate point} (e.g.\ reporting
    $[0.5,0.5]$). For a trustworthy coupled bracket, use coupled Credal~VE (outer) instead.
  \item \textbf{D5} (\texttt{CredalJT}): a separate junction-tree NLP engine, not reachable from
    ApproxLP; the reliable exact route for cross-family coupling and conditional queries.
\end{itemize}

\begin{center}\small
\begin{tabular}{lllll}
\toprule
Scheme & Where it lives for ApproxLP & Gap closed & Character & Cost \\
\midrule
D1 \texttt{method="linear-tight"} & build (\texttt{from\_lcn}) & A (partial) & free, inner tightens & $\approx 0$ \\
D2 \texttt{merge\_budget=$k$}     & build (\texttt{from\_lcn}) & A $(+\,$B$)$ & free & $2^k$ \\
D3 \texttt{solver="scip"}         & build (\texttt{from\_lcn}) & A (exact)   & free & $2^n$/solve \\
D4 \texttt{coupling="cross-family"} & inside \texttt{run} (LP rows) & B & \textbf{experimental, not certified} & med \\
D5 \texttt{CredalJT}              & separate engine            & A $+$ B     & exact, separate & treewidth \\
\bottomrule
\end{tabular}
\end{center}

\section{Verdict}\label{sec:verdict}

\begin{proposition}[Summary]\label{prop:verdict}
Let $\mathcal L$ be a chain-graph LCN and run \ALP{}.
\begin{enumerate}
  \item (\emph{Always}) With \texttt{coupling="off"}, ApproxLP returns an \emph{inner} approximation
    of $[\lo q_{\mathrm{SE}},\hi q_{\mathrm{SE}}]$ --- \emph{not} a valid outer bound on $\DLCN$; it
    can be too tight (Proposition~\ref{prop:inner}).
  \item (\emph{Tight}) When $\DSE=\DLCN$ (singly connected, no cross-family sentence) and coordinate
    descent reaches the global $\DSE$ extremum, the inner interval coincides with the exact LCN
    bounds (Proposition~\ref{prop:tight}); verified on \texttt{d4\_biting}, where ApproxLP matches the
    exact $P(B)=[0.04,0.72]$ (as does the fixed \texttt{IntervalBP}; CCTE remains inner there).
  \item (\emph{Strict gap}) The inner bound is strict exactly when the nonconvex multilinear program
    makes descent stall at a non-global local optimum; then ApproxLP under-reports the width.
  \item (\emph{Bracket}) Pairing ApproxLP (inner) with Credal~VE (exact over $\DSE$) on the same
    vertices brackets the truth and detects stalls (Section~\ref{sec:pairing}).
  \item (\emph{Tightening}) D1/D2/D3 tighten ApproxLP for free at build time; D4 is implemented but
    experimental and uncertified; D5 (\texttt{CredalJT}) is the separate exact engine.
\end{enumerate}
\end{proposition}

\noindent Placed against its siblings: Credal~VE and CCTE ($\epsilon=0$) are \emph{exact} over
$\DSE$; \texttt{IntervalBP} is \emph{outer} (interval) / \emph{inner} (variational) --- exact on
trees now that its message bugs are fixed; ApproxLP is the family's \emph{inner} member --- and, on
the instances tested here, a \emph{sound and tight} one that is the natural inner half of a two-sided
bracket. The
three companion notes cover the neighbors: \texttt{docs/cve\_exactness.tex} (the exact-over-$\DSE$
per-target engine), \texttt{docs/ibp\_exactness.tex} (the message-passing engine), and
\texttt{docs/tighter\_approximation.tex} (the D1--D5 taxonomy and the D5 exactness theorem).

\end{document}
